\documentclass[11pt]{article}\usepackage{layout}\begin{document}
\begin{canvas}\Header{第一章 · 开源与 AI}{blue}\Footer{11}
\pdfbookmark[1]{1.1　人工智能时代的开源协作：双向重塑}{page1}
\Box{15}{27}{180}{\fontsize{19.7}{25}\selectfont\color{ink}\sffamily \bfseries 1.1　人工智能时代的开源协作：双向重塑}
\Box{15}{39}{180}{\fontsize{10.7}{14.5}\selectfont\color{teal}\sffamily \bfseries AI 改变创造与协作的节奏，开源为 AI 提供可接续的基础}
\Box{15}{47.5}{180}{\fontsize{10.3}{15.5}\selectfont\color{ink}AI 正将部分编程活动从逐行编写推进到任务委托；开源则通过公开代码、工具与训练材料，让更多团队参与改进。理解这种双向关系，需要同时看工作如何完成，以及成果如何被继续使用。}
\Box{15}{66}{180}{\includegraphics[width=180mm]{.build/assets/figure-1-1.pdf}}
\Box{15}{195}{180}{\fontsize{8}{11.6}\selectfont\color{muted}\sffamily 图 1-1　AI 进入协作流程：工作入口、任务委托与责任分工。来源：GitHub Octoverse 2026\href{https://github.blog/news-insights/octoverse/octoverse-commits-pushes-reach-new-highs-in-2026-as-developers-build-with-agents/}{［S1］}及官方产品说明\href{https://docs.github.com/en/copilot/responsible-use/agents}{［S2］}。本报告主统计期为 2025.07.01—2026.06.30；右图沿用外部年度口径。流程为机制示意，不能据此推断质量或生产率的因果变化。}
\Box{15}{213}{87}{\fontsize{10.5}{15}\selectfont\color{blue}\sffamily \bfseries 1　工作入口变了，观察范围也要说清}
\Box{15}{219}{87}{\fontsize{10.2}{15.4}\selectfont\color{ink}以智能体草稿发起的 PR 比例由 3 月的 17\% 增至 9 月的超过 40\%。前者位于主统计期，后者是期后信号。这反映修改方案的发起方式，不能直接代表开源代码总量或全部开发者。}
\Box{108}{213}{87}{\fontsize{10.5}{15}\selectfont\color{blue}\sffamily \bfseries 2　从辅助回答，走向多步骤执行}
\Box{108}{219}{87}{\fontsize{10.2}{15.4}\selectfont\color{ink}在 Copilot 月活用户中，使用智能体功能的比例升至 55\%。它提示工具角色正在扩展：用户提出目标，工具组织步骤并执行。用户使用比例与 PR 起草比例分母不同，也不等同于效率提升。}
\Box{15}{250}{180}{\fontsize{10.5}{15}\selectfont\color{blue}\sffamily \bfseries 3　从生成到使用，审查与维护承接后续责任}
\Box{15}{256}{180}{\fontsize{10.2}{15.4}\selectfont\color{ink}智能体可以起草、测试并提出修改，维护者仍需核对目标、功能与安全，再决定是否采用。社区随后承担发布、文档、依赖维护与反馈处理。因而，评价 AI 时代的开源协作，应继续追踪修改是否被采纳、关系是否持续、成果是否可用；生成数量只是这条链上的一个信号。}
\end{canvas}
\begin{canvas}\Header{第一章 · 开源与 AI}{blue}\Footer{12}
\pdfbookmark[1]{开放供给，从模型延伸到基础工具}{page2}
\Box{15}{27}{180}{\fontsize{19.7}{25}\selectfont\color{ink}\sffamily \bfseries 开放供给，从模型延伸到基础工具}
\Box{15}{39}{180}{\fontsize{10.7}{14.5}\selectfont\color{teal}\sffamily \bfseries 1.1 / 开源支撑 AI　｜　同一系统中的三类工程能力}
\Box{15}{47.5}{180}{\fontsize{10.3}{15.5}\selectfont\color{ink}模型要进入实际运行环境，还需要高效计算、芯片协同与数据存取。DeepSeek 在 2025 年 2 月公开的一组基础组件，提供了具体例子：开放的不只是模型文件，还包括支撑训练与服务的工程方法。}
\Box{15}{66}{180}{\includegraphics[width=180mm]{.build/assets/figure-1-2.pdf}}
\Box{15}{195}{180}{\fontsize{8}{11.6}\selectfont\color{muted}\sffamily 图 1-2　从模型能力到系统运行：DeepSeek 五个开放基础组件的工程职责。来源：DeepSeek 官方基础设施索引“2025 年 2 月开源周”\href{https://github.com/deepseek-ai/open-infra-index\#202502-open-source-week}{［S3］}。这是主统计期之前的历史案例；按功能归类，不表示完整训练系统、全部依赖关系或性能比较。}
\Box{15}{213}{87}{\fontsize{10.5}{15}\selectfont\color{blue}\sffamily \bfseries 1　计算优化，也是一种可复用供给}
\Box{15}{219}{87}{\fontsize{10.2}{15.4}\selectfont\color{ink}FlashMLA 与 DeepGEMM 把部分底层优化方法做成可检视的代码。后续团队可以理解、测试和改进实现。能否用于自己的系统，需要核对硬件、模型与软件版本，不能仅凭公开代码判断部署成本。}
\Box{108}{213}{87}{\fontsize{10.5}{15}\selectfont\color{blue}\sffamily \bfseries 2　多芯片能力，依赖协同而非简单叠加}
\Box{108}{219}{87}{\fontsize{10.2}{15.4}\selectfont\color{ink}DeepEP 关注芯片间的数据交换，DualPipe 关注训练中计算与通信的安排。二者对应不同瓶颈。开放这些实现，有助于其他团队沿既有方法开展适配；实际收益仍取决于网络、任务和工程条件。}
\Box{15}{250}{180}{\fontsize{10.5}{15}\selectfont\color{blue}\sffamily \bfseries 3　可持续服务，需要看见数据与维护环节}
\Box{15}{256}{180}{\fontsize{10.2}{15.4}\selectfont\color{ink}3FS 支撑数据装载、训练进度保存与恢复，以及推理中的缓存存取。它提醒我们：模型表现之外，系统能否持续运行还取决于数据和基础设施。评估开放供给，需要进一步观察采用者能否复现环境、完成适配并持续维护。这也是第四章从技术体系而非单个项目理解生态的原因。}
\end{canvas}
\begin{canvas}\Header{第一章 · 开源与 AI}{blue}\Footer{13}
\pdfbookmark[1]{中国开放模型，进入全球再创造}{page3}
\Box{15}{27}{180}{\fontsize{19.7}{25}\selectfont\color{ink}\sffamily \bfseries 中国开放模型，进入全球再创造}
\Box{15}{39}{180}{\fontsize{10.7}{14.5}\selectfont\color{teal}\sffamily \bfseries 1.1 / 跨组织复用　｜　从生成回答到新的训练材料}
\Box{15}{47.5}{180}{\fontsize{10.3}{15.5}\selectfont\color{ink}模型的输出可以被整理为另一模型的训练材料。NVIDIA 于 2025 年 5 月公开说明：Qwen、DeepSeek-R1 等模型生成的记录，经整理进入 Llama-Nemotron 后训练数据集，并用于后续模型训练。}
\Box{15}{66}{180}{\includegraphics[width=180mm]{.build/assets/figure-1-3.pdf}}
\Box{15}{195}{180}{\fontsize{8}{11.6}\selectfont\color{muted}\sffamily 图 1-3　开放模型输出的再利用：生成来源、样本记录与后续成果。来源：NVIDIA 技术说明（2025-05-14，表 2）\href{https://developer.nvidia.com/blog/build-custom-reasoning-models-with-advanced-open-post-training-datasets/}{［S4］}及历史数据集卡\href{https://huggingface.co/datasets/nvidia/Llama-Nemotron-Post-Training-Dataset/blob/04d085443485c2b8e14ad5f518005c769deb4904/README.md}{［S5］}。样本数选取三个生成模型的原表条目；这是历史复用案例，不作国别、市场或性能贡献份额推导。}
\Box{15}{213}{87}{\fontsize{10.5}{15}\selectfont\color{blue}\sffamily \bfseries 1　中国贡献，可以沿具体资源追踪}
\Box{15}{219}{87}{\fontsize{10.2}{15.4}\selectfont\color{ink}这个案例把中国模型与海外团队的后续研发联系起来。Qwen、DeepSeek-R1 与其他来源共同提供生成能力，NVIDIA 继续整理材料并训练模型。跨组织接续由可查的资源和使用关系体现。}
\Box{108}{213}{87}{\fontsize{10.5}{15}\selectfont\color{blue}\sffamily \bfseries 2　样本规模，说明供给而非效果分摊}
\Box{108}{219}{87}{\fontsize{10.2}{15.4}\selectfont\color{ink}图中的一条样本是一条训练记录；同一问题可以产生不同回答。约 121.3 万条 R1 样本说明其输出进入了后续训练过程。数量不能直接换算为模型能力提升，也不能作为整个生态的贡献比例。}
\Box{15}{250}{180}{\fontsize{10.5}{15}\selectfont\color{blue}\sffamily \bfseries 3　开放材料越完整，越有利于下一次接续}
\Box{15}{256}{180}{\fontsize{10.2}{15.4}\selectfont\color{ink}公开样本、生成与整理工具及后续模型，为其他团队复现方法和继续改进提供了条件。报告由此沿“参与、供给、复用”追踪资源流动，并继续检验治理与公共价值。要让开放成果成为可持续公共能力，还需解决接入成本、本地适配与长期维护；这构成下一节讨论发展目标的起点。}
\end{canvas}
\end{document}